\section{Sprint-by-Sprint Facilitation Notes}
\textit{Reference $|$ Course Text: Chapters referenced per sprint below}

This section provides a half-page summary for each sprint: what to look for, red flags, key deliverables, and suggested questions for sprint reviews. Use this as your pre-meeting preparation checklist.

% ── SDI ─────────────────────────────────────────────────────────────

\subsection{Sprint 0: Project Expo, Team Formation, Client Kickoff}
\textit{Weeks 1--2 $|$ Course Text: Chapters 1, 2}

\textbf{Health indicators:} Team has contacted the client. Meeting logistics are confirmed. Team members have introduced themselves and begun discussing roles.

\begin{redflag}
Team has not contacted the client by end of Week~1. A team that cannot send an email will struggle with every subsequent communication milestone.
\end{redflag}

\textbf{Key deliverables:} Project Selection Survey, client introduction email with team slide, initial meeting scheduled.

\begin{paquestion}
\begin{itemize}[nosep,leftmargin=1.5em]
\item What does the client actually need---not what do they say they want?
\item What are the top three unknowns about this project right now?
\item How will you prepare for the kickoff meeting?
\end{itemize}
\end{paquestion}

% ──────────────────────────────────────────────────────

\subsection{Sprint 1: Client Discovery, Problem Framing}
\textit{Weeks 3--4 $|$ Course Text: Chapters 2, 3}

\textbf{Health indicators:} Kickoff meeting completed. Meeting minutes distributed. Team can articulate the problem in their own words (not just echoing the client brief).

\begin{redflag}
Team is already proposing solutions. Sprint~1 is for understanding the problem, not solving it. If you hear ``we're going to use an Arduino,'' redirect to ``what problem are you solving?''
\end{redflag}

\textbf{Key deliverables:} Kickoff meeting minutes, initial needs/goals documentation, stakeholder identification.

\begin{paquestion}
\begin{itemize}[nosep,leftmargin=1.5em]
\item Who are the stakeholders beyond the primary client contact?
\item What constraints has the client stated explicitly? What constraints are implied?
\item What does success look like from the client's perspective?
\end{itemize}
\end{paquestion}

% ──────────────────────────────────────────────────────

\subsection{Sprint 2: Requirements Development, Research}
\textit{Weeks 5--6 $|$ Course Text: Chapters 3, 4}

\textbf{Health indicators:} Requirements are written in ``shall'' statement form with measurable acceptance criteria. The team distinguishes between requirements and specifications. Background research is underway.

\begin{redflag}
Requirements that read like specifications (``The system shall use a 12V motor'') instead of functional needs (``The system shall generate a minimum torque of 5~N$\cdot$m''). This is the single most common Sprint~2 error.
\end{redflag}

\textbf{Key deliverables:} Draft SDRD (System Design Requirements Document), literature review progress, concept generation begins.

\begin{paquestion}
\begin{itemize}[nosep,leftmargin=1.5em]
\item How will you verify each requirement? If you cannot test it, rewrite it.
\item What engineering standards apply to this project?
\item What did your background research reveal that surprised you?
\end{itemize}
\end{paquestion}

% ──────────────────────────────────────────────────────

\subsection{Sprint 3: Concept Development, Client Approval}
\textit{Weeks 7--8 $|$ Course Text: Chapters 3, 4, 17}

\textbf{Health indicators:} Multiple concepts generated (target $\geq$15 ideas before down-selection). Decision matrix uses weighted criteria with sensitivity analysis. Client has reviewed and approved the selected approach.

\begin{redflag}
Only one concept considered, or the decision matrix is clearly rigged to justify a predetermined choice. Ask the team to vary the weights by $\pm$20\% and report whether the winner changes.
\end{redflag}

\textbf{Key deliverables:} Decision matrix with sensitivity analysis, client approval of approach, engineering standards identification.

\begin{paquestion}
\begin{itemize}[nosep,leftmargin=1.5em]
\item What were the top three concepts you eliminated, and why?
\item If you change the top-weighted criterion, does your winner change?
\item What is the client's biggest concern about the selected concept?
\end{itemize}
\end{paquestion}

% ──────────────────────────────────────────────────────

\subsection{Sprint 4: Preliminary Design, PoC Planning}
\textit{Weeks 9--10 $|$ Course Text: Chapters 4, 6, 7}

\textbf{Health indicators:} SOW is drafted with clear scope boundaries (including an ``Out of Scope'' section). The Proof of Concept (PoC) targets the project's key technical risk, not the easiest thing to build.

\begin{redflag}
PoC plan that tests something the team already knows works. The purpose of a PoC is to retire the \emph{highest} technical risk. If the team's PoC does not make them nervous, it is probably targeting the wrong risk.
\end{redflag}

\textbf{Key deliverables:} Draft SOW, PoC plan, initial project schedule (Gantt or backlog).

\begin{paquestion}
\begin{itemize}[nosep,leftmargin=1.5em]
\item What is the single biggest technical risk in this project?
\item Does your PoC directly address that risk?
\item What is explicitly out of scope, and has the client agreed?
\end{itemize}
\end{paquestion}

% ──────────────────────────────────────────────────────

\subsection{Sprint 5: PDR Preparation, PoC Execution}
\textit{Weeks 11--12 $|$ Course Text: Chapters 7, 8}

\textbf{Health indicators:} PoC is built and tested (not just planned). PDR report draft is under internal review. The team has scheduled a mock PDR.

\begin{redflag}
Team skips the mock PDR. Mock PDRs consistently improve final PDR quality. If a team claims they ``don't have time,'' they especially need one.
\end{redflag}

\textbf{Key deliverables:} PoC results, PDR report draft, PDR presentation draft, mock PDR completed.

\begin{paquestion}
\begin{itemize}[nosep,leftmargin=1.5em]
\item What did the PoC teach you? Did anything fail, and what did you learn from the failure?
\item Walk me through your PDR presentation flow in 60 seconds.
\item What question are you most afraid a reviewer will ask at PDR?
\end{itemize}
\end{paquestion}

% ──────────────────────────────────────────────────────

\subsection{Sprint 6: PDR, Long-Lead Procurement, Semester Wrap}
\textit{Weeks 13--14 $|$ Course Text: Chapters 8, 14, 15}

\textbf{Health indicators:} PDR is complete. Feedback has been documented and an action plan exists. Long-lead items are ordered or have purchase orders submitted. Budget is realistic.

\begin{redflag}
No procurement activity after PDR. If the team has not ordered long-lead components (PCBs, custom parts, specialty materials) by end of Sprint~6, those items will not arrive in time for SDII fabrication.
\end{redflag}

\textbf{Key deliverables:} PDR presentation and report, PDR feedback action plan, procurement initiation, semester reflection.

\begin{paquestion}
\begin{itemize}[nosep,leftmargin=1.5em]
\item What were the top three weaknesses identified at PDR, and what is your plan to address each?
\item What is your longest lead-time item, and when will it arrive?
\item What will you accomplish over the break to stay on track?
\end{itemize}
\end{paquestion}

\begin{paaction}
Complete Peer Evaluation \#2 and Individual Performance Evaluation \#1 before the semester ends. Submit your PA grade input to the CF by the posted deadline.
\end{paaction}

% ── SDII ────────────────────────────────────────────────────────────

\subsection{Sprint 7: SDII Kickoff, Detailed Design, Hazard Assessment}
\textit{Weeks 15--16 $|$ Course Text: Chapters 9, 12, 21}

\textbf{Health indicators:} Team has addressed PDR feedback. Detailed design is progressing (CAD, schematics, analysis). Hazard Assessment is drafted. Fabrication plan is realistic.

\begin{redflag}
Team wants to start fabrication without a completed Hazard Assessment. This is a hard stop. No fabrication may begin until the Hazard Assessment is reviewed and approved by the PA.
\end{redflag}

\textbf{Key deliverables:} PDR feedback closure, detailed design progress, Hazard Assessment, updated project schedule.

\begin{paquestion}
\begin{itemize}[nosep,leftmargin=1.5em]
\item Which PDR feedback items have you closed? Which are still open?
\item Walk me through your Hazard Assessment. What is the highest-severity hazard?
\item Is your fabrication timeline realistic given InnoHub availability and material lead times?
\end{itemize}
\end{paquestion}

% ──────────────────────────────────────────────────────

\subsection{Sprint 8: Fabrication Begins, CDR Preparation}
\textit{Weeks 17--18 $|$ Course Text: Chapters 9, 12, 14}

\textbf{Health indicators:} Fabrication is underway (not just planned). BOM is complete---including fasteners, consumables, and wiring. CDR report draft is in progress.

\begin{redflag}
BOM missing fasteners and consumables. Have the team walk through the CAD assembly and identify every component that needs a BOM line. If it is visible in the model, it needs a part number and a cost.
\end{redflag}

\textbf{Key deliverables:} Fabrication progress, complete BOM, CDR report draft, CDR presentation draft.

\begin{paquestion}
\begin{itemize}[nosep,leftmargin=1.5em]
\item Show me your BOM. Is every component accounted for, including fasteners?
\item What is your integration plan---how will subsystems come together?
\item What is your budget status? Any overruns or pending purchases?
\end{itemize}
\end{paquestion}

% ──────────────────────────────────────────────────────

\subsection{Sprints 9--11: Build, Integration, Testing}
\textit{Weeks 19--24 $|$ Course Text: Chapters 6, 9, 12}

\textbf{Health indicators:} Subsystem testing is occurring in parallel with fabrication. Integration is incremental (one subsystem at a time), not big-bang. V\&V testing has started---not deferred to Sprint~12.

\begin{redflag}
All testing deferred to ``after we finish building.'' Testing should begin as soon as the first subsystem is functional. Every week of testing margin is one redesign cycle. Teams that leave all testing to Sprint~12 run out of time to fix problems.
\end{redflag}

\textbf{Key deliverables:} Subsystem test results, integration progress, V\&V test execution, updated VCRM.

\begin{paquestion}
\begin{itemize}[nosep,leftmargin=1.5em]
\item Which subsystems have been tested independently? What were the results?
\item What is your integration sequence, and what are the interface risks?
\item How many requirements have you verified so far? How many remain?
\end{itemize}
\end{paquestion}

\begin{paaction}
During these sprints, shift your focus from design guidance to project management oversight. Ask to see the Gantt chart or backlog at every meeting. Teams that lose schedule visibility in Sprints 9--11 are the ones that scramble at FDR.
\end{paaction}

% ──────────────────────────────────────────────────────

\subsection{Sprints 12--13: Final Testing, Documentation, FDR Prep}
\textit{Weeks 25--28 $|$ Course Text: Chapters 6, 10, 18}

\textbf{Health indicators:} V\&V testing is substantially complete. VCRM shows all requirements dispositioned. FDR report sections are being written concurrently with testing (not deferred to the last week). Documentation debt is minimal.

\begin{redflag}
Documentation debt is the number one indicator of a project in trouble at this stage. If the team has been building for six sprints without writing report sections, they will produce a rushed, low-quality FDR report. Check documentation status at every meeting.
\end{redflag}

\textbf{Key deliverables:} Completed V\&V results, VCRM with all requirements dispositioned, FDR report draft, FDR presentation draft.

\begin{paquestion}
\begin{itemize}[nosep,leftmargin=1.5em]
\item Show me the VCRM. Which requirements are verified, failed, waived, or deferred?
\item For any failed requirements, what is the disposition and path forward?
\item How many pages of the FDR report are drafted? How many remain?
\end{itemize}
\end{paquestion}

% ──────────────────────────────────────────────────────

\subsection{Sprint 14: FDR, Design Showcase, Course Closeout}
\textit{Weeks 29--30 $|$ Course Text: Chapters 10, 19}

\textbf{Health indicators:} FDR report is complete and sent to the client. Presentation is rehearsed. Demo hardware is tested and working. The team can articulate what they would do differently.

\begin{redflag}
Incomplete VCRM at FDR. Every requirement must be dispositioned: Verified, Failed (with justification and recommendation), Waived (with client agreement), or Deferred (with rationale). An incomplete VCRM signals that the team did not complete the design cycle.
\end{redflag}

\textbf{Key deliverables:} FDR presentation and report, Design Showcase poster and demo, final client handoff, course evaluations.

\begin{paquestion}
\begin{itemize}[nosep,leftmargin=1.5em]
\item If you had one more sprint, what would you change?
\item What was the most significant technical challenge, and how did you resolve it?
\item What advice would you give to next year's teams?
\end{itemize}
\end{paquestion}

\begin{paaction}
Complete all remaining Peer Evaluations and Individual Performance Evaluations. Submit final PA grade input to the CF by the posted deadline. Ensure client handoff documentation is complete.
\end{paaction}


% ═══════════════════════════════════════
